\chapter{Related Work}
\label{ch:related}

This chapter reviews five areas that are directly related to this thesis: psychiatric LLM diagnosis, retrieval with similar examples, medical multi-agent systems, criterion-grounded audit records, and primary-diagnosis selection. The goal is not to rank the reviewed systems. The goal is to show which parts of the diagnostic process they evaluate and which questions remain open.

\section{Psychiatric LLM Diagnosis}

Psychiatric LLM studies cover several different tasks. Some test psychiatric knowledge or licensing questions. Others use short clinical vignettes, focus on one disorder, detect depression from text, diagnose from a fixed record, or support an interactive consultation~\citep{qiu2023largeai,li2024taiwan,raballo2025,sarma2025,li2025addreview,mcduff2025amie}. These tasks should not be treated as the same problem. They differ in the evidence provided to the model, the number of possible diagnoses, the required output, and the reference used for scoring.

Screening is also different from fine-grained differential diagnosis. Detecting a strong depression or anxiety signal is easier than choosing one primary diagnosis among several disorders that share symptoms. A disorder-level decision may depend on symptom duration, which symptoms began first, illness course, exclusion conditions, and comorbidity. LingxiDiagBench shows this difference: performance falls as the task moves from broad depression--anxiety classification to comorbidity recognition and twelve-category diagnosis~\citep{xu2026lingxidiagbench}. MDD-5k provides another synthetic Chinese dialogue resource generated from anonymized patient cases~\citep{yin2024mdd5k}. Both datasets support repeatable experiments, but neither replaces evaluation on natural clinical conversations.

Another important difference is whether the system can ask for more information. MIND, WiseMind/ProAI, MAGI, DSM5AgentFlow, and conversational diagnostic AI study interactive workflows that can ask questions during the assessment~\citep{li2026mind,wu2026wisemind,bi2025magi,ozgun2025dsm5agentflow,tu2025conversational}. MentalSeek-Dx and MoodAngels instead focus on information that is already available in a record or transcript~\citep{sun2026mentalseek,xiao2025moodangels}. Interactive systems combine evidence collection with diagnosis. Fixed-input systems must work with the same possibly incomplete record. The fixed-input setting used in this thesis therefore tests how the model uses the available transcript without giving it credit for asking additional questions.

Most published evaluations still focus on the final diagnosis. They provide less information about whether a benchmark disagreement began because a diagnosis was never proposed, did not pass a criterion rule, or remained available but was not selected as primary. This gap motivates a stage-wise evaluation rather than final-label accuracy alone.

\section{Retrieval-Augmented Diagnostic Reasoning}

Retrieval can play two different roles. Knowledge retrieval brings in documents, guidelines, or other factual sources~\citep{lewis2020rag}. Demonstration retrieval selects labeled examples or similar cases and places them in the prompt for in-context learning~\citep{liu2022incontext}. This thesis uses the second role. Retrieved cases are examples from the training source, not medical evidence about the current patient.

Similar examples can show common wording, expected output format, possible labels, and frequent symptom--diagnosis patterns. However, in-context learning depends strongly on which examples are selected~\citep{liu2022incontext}. A retrieved case describes another patient or synthetic case. It may guide the model, but it should not be treated as evidence that the current patient has the same diagnosis.

Retrieval is not always helpful. The retrieved examples reflect the wording and label distribution of the indexed corpus. Frequent classes may appear many times. Strong lexical overlap may encourage shortcuts, and a weakly related example may distract the model. Retrieval should therefore be tested as a configuration rather than assumed to improve diagnosis.

A label-aware policy may show a wider range of diagnoses and reduce repeated examples from frequent classes. However, balancing changes more than label frequency. It may also change the number, order, similarity, total length, and mix of examples. A comparison between global similarity retrieval and a balanced policy is therefore a comparison between complete retrieval settings, not a clean estimate of the effect of label balance alone.

Prior work shows that example choice matters for in-context learning, but it does not establish which retrieval policy works best for fine-grained psychiatric diagnosis. Retrieval policies should therefore be compared within the same architecture and evaluation contract.

\section{Medical Multi-Agent Reasoning}

Medical multi-agent systems organize several model calls in different ways. Some assign separate roles for diagnosis, evidence extraction, criticism, or final decision-making. Others use consultation, debate, voting, or consensus. Verification-oriented systems produce an answer and then ask another component to check or revise it. MedAgents uses several role-based medical experts, and MDAgents changes the collaboration structure according to estimated task difficulty~\citep{tang2024medagents,kim2024mdagents}. Multi-agent debate is a related general method in which agents exchange and revise competing answers~\citep{du2024debate}.

These designs can be useful. A narrow role may focus on one subtask. A second model call may find missing evidence. Cross-checking may compare a diagnosis with criteria, and a structured workflow can save intermediate outputs for review. These are possible benefits, not guarantees. More agents or more discussion rounds do not automatically produce better diagnostic accuracy.

Multi-agent systems can also repeat the same mistakes. Agents that use the same backbone model, similar prompts, and the same evidence may have correlated errors. A wrong candidate or unsupported evidence statement can pass into later stages. Free-form discussion may change several parts of the workflow at once, and the final component may still choose the wrong answer even when a useful alternative was mentioned earlier.

This creates an attribution problem. When the candidate generator, evidence representation, discussion process, and final aggregation all change together, a difference in final accuracy is a complete-system effect. It does not show which component caused the difference. A comparison with a Single LLM may also change prompt structure, context, model-call count, and compute.

For this reason, candidate generation, criterion checking, and final commitment should be recorded separately. This does not make the components statistically independent, but it allows the observed disagreement to be located more precisely.

\section{Auditability and Criterion Grounding}
\label{sec:auditability-definition}

A fluent explanation is not proof that the explanation reflects the factors that produced the answer. Model-generated reasoning can sound clear while still being unfaithful~\citep{turpin2023unfaithful}. In this thesis, \emph{explainability} means a human-facing account of a decision. \emph{Interpretability} means how directly a person can understand the decision process. \emph{Faithfulness} asks whether the stated account reflects the factors that produced the output. \emph{Auditability} means that an external record is saved and can be inspected or replayed. These properties are different and should be evaluated separately in healthcare~\citep{rudin2019stop,huang2024xiai}.

Diagnostic-chain and symptom-based methods can show intermediate hypotheses, confidence values, or clinically motivated features~\citep{chen2025cod,lee2026interpretable,lyu2026depllm}. Psychiatric systems may also connect outputs to guidelines, structured interviews, symptom-to-criterion mappings, and supporting or conflicting evidence~\citep{li2026mind,bi2025magi,ozgun2025dsm5agentflow}. Constraint-logic methods make the rule layer more visible by translating diagnostic criteria into inspectable programs~\citep{kim2025clp}.

These records are useful only if the underlying criteria and evidence mapping are reasonable. A model may generate an unsupported evidence note. A criterion may be represented incorrectly, and an inspectable rule may still contain a poor clinical assumption. Expert review of model-generated diagnostic rules remains necessary~\citep{kim2025clp}. Criterion-grounded outputs are therefore reviewable model records, not clinician-validated evidence or clinical truth.

An auditable system should keep the source evidence, criterion states, rule results, intermediate diagnoses, and final decision. This record still does not reveal hidden model reasoning. A criterion layer may be used only for review, or it may take part in final selection. These two roles should be evaluated separately because good evidence recording or high compatibility inclusion does not guarantee a good primary diagnosis.

The remaining research gap is therefore not only whether intermediate records exist. It is whether candidate availability, criterion compatibility, and final commitment are evaluated together on the same cases.

\section{Primary-Diagnosis Selection}

Criterion compatibility and primary-diagnosis selection answer different questions. Compatibility asks whether one diagnosis fits the available transcript under a defined rule. Because several psychiatric disorders can share symptoms, more than one diagnosis may pass its own rule. Compatibility alone does not show whether a diagnosis should be primary, comorbid, secondary, or only a differential alternative.

Structured medical knowledge can help organize evidence. MedKGI combines a medical knowledge graph, guided questioning, and a structured dialogue state to seek useful information during diagnosis~\citep{wang2025medkgi}. DKEC builds heterogeneous knowledge graphs from external medical sources and uses their relations for multilabel diagnosis prediction~\citep{ge2024dkec}. These studies show how structured representations can support evidence collection or label prediction. Their evaluated goals are different from selecting one primary diagnosis among several criterion-compatible psychiatric candidates in a fixed transcript.

A primary selector may need direct comparisons among candidates. Important relations include symptom timing, duration, illness course, previous episodes, exclusionary causes, syndrome dominance, explanatory dependence, and comorbidity. The reviewed methods can organize evidence or identify plausible labels, but they do not directly test the case in which a benchmark-reference diagnosis remains both ranked and criterion-compatible while another diagnosis is committed as primary.

This gap motivates separate measurement of genuine Top-3 candidate coverage, benchmark-reference inclusion in the criterion-compatible set, and committed-primary agreement.

\section{Research Gap and Research Questions}

The reviewed work exposes different parts of the diagnostic process. Some systems produce hypotheses, some record criterion-related evidence, and some return only a final diagnosis. Table~\ref{tab:related-work-matrix} summarizes the outputs reported in representative studies. The table describes published evaluation contracts. It does not rank overall system quality or clinical ability.

\begin{table}[htbp]
\centering
\caption{Representative psychiatric diagnostic systems and the outputs available for stage-wise disagreement analysis.}
\label{tab:related-work-matrix}
\small
\begin{adjustbox}{max width=\textwidth}
\begin{tabular}{|>{\raggedright\arraybackslash}p{2.5cm}|>{\raggedright\arraybackslash}p{2.7cm}|>{\raggedright\arraybackslash}p{3.1cm}|>{\raggedright\arraybackslash}p{4.3cm}|>{\centering\arraybackslash}p{2.0cm}|}
\hline
Study or system & Evidence setting & Ranked candidates & Criterion or evidence representation & Stage-wise evaluation\\
\hline
MIND~\citep{li2026mind} & Interactive consultation & No ranked evaluation reported & Criteria-grounded rationales and process rewards & No\\
\hline
WiseMind/\allowbreak ProAI~\citep{wu2026wisemind} & Interactive consultation & No ranked evaluation reported & Structured DSM-guided action states & Partial\\
\hline
MAGI~\citep{bi2025magi} & Interactive MINI-guided interview & No ranked evaluation reported & Symptom and dialogue evidence mapped to criteria & Partial\\
\hline
DSM5Agent\allowbreak Flow~\citep{ozgun2025dsm5agentflow} & Interactive questionnaire & No ranked evaluation reported & Utterance-level supporting and contradictory criteria & Partial\\
\hline
MentalSeek-Dx~\citep{sun2026mentalseek} & Fixed clinical record & Hypotheses generated; ranking not evaluated & Criteria-guided staged reasoning & Partial\\
\hline
Mood\allowbreak Angels~\citep{xiao2025moodangels} & Fixed records and assessment scales & No ranked evaluation reported & DSM-5 retrieval and symptom matching & Partial\\
\hline
Constraint-logic CDSS~\citep{kim2025clp} & Structured patient facts & None & Inspectable diagnostic rules & No\\
\hline
\end{tabular}
\end{adjustbox}

\vspace{0.4em}
\begin{minipage}{0.98\textwidth}
\footnotesize
\textit{Note.} ``Partial'' means that a study reports intermediate diagnostic or criterion-related outputs but does not separately evaluate candidate detection, criterion checking, and primary commitment on the same cases. ``Not reported'' refers only to the published evaluation and does not show that the implementation lacks the capability.
\end{minipage}
\end{table}
\FloatBarrier

The reviewed literature leaves five questions that are directly related to this thesis:

\begin{enumerate}
    \item Complete-system comparisons are difficult because conventional classifiers, Single LLMs, and multi-stage systems often use different prompts, model calls, retrieval settings, and output contracts.
    \item Final-label accuracy does not show whether a reference diagnosis was missing from the candidates, absent from the criterion-compatible set, or available in both views but not committed as primary.
    \item Demonstration retrieval depends on which examples are selected, so global similarity retrieval and label-aware balancing should be compared within each architecture.
    \item Multi-agent and criterion-based selection methods often change several parts of the workflow at once, which limits causal claims about one final-selection component.
    \item A stage-wise analysis should be checked on a second synthetic source because one corpus cannot show whether the same output views remain useful after the data source changes.
\end{enumerate}

HiED addresses these questions by saving the ranked candidates, criterion states, criterion-compatible set, and final commitment under declared comparison contracts.

This thesis uses the following research questions:

\begin{description}[style=nextline,leftmargin=1.5cm,labelwidth=1.1cm,itemsep=0pt,parsep=0pt,topsep=0pt,partopsep=0pt]
\item[RQ1] How does HiED perform relative to conventional classification and a Single LLM on the same internal split and parent-label evaluation?
\item[RQ2] Where is benchmark disagreement recorded across candidate generation, criterion checking, compatibility analysis, and primary commitment?
\item[RQ3] How do similar-case retrieval strategies affect Single and HiED?
\item[RQ4] Do the tested primary-selection strategies reduce the recorded selection disagreement?
\item[RQ5] Does the stage-wise analysis remain useful under a second synthetic source?
\end{description}

RQ3 first compares retrieval settings on the public-validation split and freezes the selected settings. RQ1 then evaluates those selected settings on the fixed internal held-out split. RQ2 uses the internal HiED trace to locate recorded disagreement. RQ4 compares the tested primary-selection methods. RQ5 applies the same output views to MDD-5k and uses the cross-corpus analyses to define the synthetic-source boundary.
